\documentclass[14pt]{extreport}
\usepackage{gost}
\usepackage{emptypage}
\usepackage{lscape}
\usepackage{ulem}
\usepackage{amssymb, amsmath}
\usepackage{graphicx}
\usepackage{url}
\usepackage{textcomp}
\usepackage{float}

\setcounter{page}{1}
\thispagestyle{empty}

\begin{document}

% ===== ТИТУЛЬНЫЙ ЛИСТ =====
\thispagestyle{empty}
\begingroup
\singlespacing
\setlength{\parskip}{0pt}

\begin{center}
    {\bfseries Министерство науки и высшего образования Российской Федерации} \\[1pt]
    ФЕДЕРАЛЬНОЕ ГОСУДАРСТВЕННОЕ АВТОНОМНОЕ ОБРАЗОВАТЕЛЬНОЕ УЧРЕЖДЕНИЕ ВЫСШЕГО ОБРАЗОВАНИЯ \\[1pt]
    {\bfseries НАЦИОНАЛЬНЫЙ ИССЛЕДОВАТЕЛЬСКИЙ УНИВЕРСИТЕТ ИТМО} \\[1pt]
    {\bfseries ITMO University} \\[0.8cm]
    {\bfseries ОТЧЕТ ПО ПРАКТИЧЕСКОЙ РАБОТЕ № 3}
\end{center}

\vspace{0.5cm}

% ---------- Блок данных (с увеличенным межстрочным интервалом) ----------
{\setstretch{1.5}
\begin{flushleft}
    {\bfseries По дисциплине:} Контейнеризация и оркестрация приложений \par
    {\bfseries Тема работы:} Персистентность и изоляция \par
    {\bfseries Обучающийся:} Лясин Илья Игоревич \par
    {\bfseries Факультет:} 11.03.02 Прикладная информатика \par
    {\bfseries Группа:} К3321 \par
    {\bfseries Направление подготовки:} 11.03.02 Инфокоммуникационные технологии и системы связи \par
    {\bfseries Образовательная программа:} Программирование в инфокоммуникационных системах \par
\end{flushleft}
}

\vspace{0.5cm}

% ---------- Блок подписей на всю ширину страницы ----------
\noindent
\begin{tabular*}{\textwidth}{@{}p{2.8cm} @{\extracolsep{\fill}} p{3.2cm} p{3.2cm} p{3.2cm}@{}}
\textbf{Обучающийся} & \makebox[3.2cm]{\hrulefill} & \makebox[3.2cm]{\hrulefill} & \makebox[3.2cm]{\shortstack{Лясин И.И.\\[-2pt]\rule{3.2cm}{0.4pt}}} \\[2pt]
                     & \makebox[3.2cm]{\scriptsize(дата)} & \makebox[3.2cm]{\scriptsize(подпись)} & \makebox[3.2cm]{\scriptsize(Ф.И.О.)} \\[1.2cm]
\textbf{Руководитель} & \makebox[3.2cm]{\hrulefill} & \makebox[3.2cm]{\hrulefill} & \makebox[3.2cm]{\shortstack{Аминов Н.С.\\[-2pt]\rule{3.2cm}{0.4pt}}} \\[2pt]
                      & \makebox[3.2cm]{\scriptsize(дата)} & \makebox[3.2cm]{\scriptsize(подпись)} & \makebox[3.2cm]{\scriptsize(Ф.И.О.)} \\
\end{tabular*}

\vfill

\begin{center}
Санкт-Петербург 2026
\end{center}
\endgroup


% ===== ОГЛАВЛЕНИЕ =====
\newpage
\tableofcontents
\thispagestyle{fancyplain}
\newpage

% ===== ВВЕДЕНИЕ =====
\intro\label{intro}
\textbf{Цель работы:} Подключаем к приложению PostgreSQL, учимся хранить данные вне контейнера и разбираемся с сетями Docker.

В рамках работы рассматриваются основные механизмы контейнеризации: именованные тома для хранения данных, bind mount для передачи конфигурации, а также пользовательские сети для изоляции и связывания контейнеров. Практическая часть выполняется с использованием персональных данных, приведенных в таблице \ref{tab:params}.

\begin{table}[H]
\caption{Персональные параметры}
\label{tab:params}
\begin{tabular}{|l|l|}
\hline
Параметр & Значение \\
\hline
Номер ИСУ & 466570 \\
Персональный порт (app) & 5500 \\
Персональный порт (assistant) & 8070 \\
Персональный префикс & lyasin \\
Секретный ключ & secret70lyasin \\
Имя базы данных & lyasin70 \\
\hline
\end{tabular}
\end{table}

% ===== ГЛАВЫ =====
\chapter{Основная часть}

\section{Этап 1. Запуск контейнера с именованным volume}

В ранее склонированном репозитории перейдём в ветку \texttt{lab-2} и создадим новую \texttt{lab-3}. Далее создадим новый том.

\begin{figure}[H]
    \centering
    \includegraphics[width=0.8\textwidth]{image-1.png}
    \caption{Создание нового тома}
    \label{fig:stage1-volume}
\end{figure}

И запустим postgres с созданным томом (рисунок~\ref{fig:stage1-postgres}). Том будет подключаться как директория \texttt{var/lib/postgresql/data} внутри контейнера. Именно эта директория содержит базы данных в postgres.

\begin{figure}[H]
    \centering
    \includegraphics[width=0.8\textwidth]{image-2.png}
    \caption{Запуск postgres с созданным томом}
    \label{fig:stage1-postgres}
\end{figure}

Запустим контейнер с приложением ассистент (рисунок~\ref{fig:stage1-assistant}). Мы передаем переменную \texttt{DATABASE\_URL} как адрес \texttt{host.docker.internal} - DNS имя хоста. Т.е. приложение получает доступ к БД через сеть хоста, также как пользователь подключается к БД через localhost.

\begin{figure}[H]
    \centering
    \includegraphics[width=0.8\textwidth]{image-3.png}
    \caption{Запуск контейнера с приложением ассистент}
    \label{fig:stage1-assistant}
\end{figure}

И проверим его работоспособность (рисунок~\ref{fig:stage1-check}).

\begin{figure}[H]
    \centering
    \includegraphics[width=0.8\textwidth]{image-4.png}
    \caption{Проверка работоспособности приложения}
    \label{fig:stage1-check}
\end{figure}

\section{Этап 2. Тест персистентности: пересоздание контейнера}

Проверим, что приложение сохраняет информацию после перезапуска контейнера (рисунок~\ref{fig:stage2-restart}).

\begin{figure}[H]
    \centering
    \includegraphics[width=0.8\textwidth]{image-5.png}
    \caption{Проверка сохранения информации после перезапуска}
    \label{fig:stage2-restart}
\end{figure}

Выполним запрос к приложению (рисунок~\ref{fig:stage2-request}).

\begin{figure}[H]
    \centering
    \includegraphics[width=0.8\textwidth]{image-6.png}
    \caption{Запрос к приложению}
    \label{fig:stage2-request}
\end{figure}

Артефакт персистентности (рисунок~\ref{fig:stage2-artifact}).

\begin{figure}[H]
    \centering
    \includegraphics[width=0.8\textwidth]{image-7.png}
    \caption{Артефакт персистентности}
    \label{fig:stage2-artifact}
\end{figure}

\section{Этап 3. Bind mount для конфигурации}

Изменим конфигурационный файл \texttt{config/app.conf} (рисунок~\ref{fig:stage3-config}).

\begin{figure}[H]
    \centering
    \includegraphics[width=0.8\textwidth]{image-8.png}
    \caption{Изменение конфигурационного файла}
    \label{fig:stage3-config}
\end{figure}

Теперь запустим новый контейнер, примонтировав конфигурацию как \texttt{bind mount}. Затем зайдём внутрь контейнера и проверим, что конфигурация действительно доступна изнутри контейнера (рисунок~\ref{fig:stage3-mount}). Конфиг подключен как read-only, т.е. конфигурация производится только извне контейнера.

\begin{figure}[H]
    \centering
    \includegraphics[width=0.8\textwidth]{image-9.png}
    \caption{Проверка доступности конфигурации внутри контейнера}
    \label{fig:stage3-mount}
\end{figure}

\section{Этап 4. Пользовательские сети для изоляции}

Создадим пользовательскую сеть и удалим ненужные контейнеры с прошлых этапов (рисунок~\ref{fig:stage4-network}).

\begin{figure}[H]
    \centering
    \includegraphics[width=0.8\textwidth]{image-10.png}
    \caption{Создание пользовательской сети и удаление контейнеров}
    \label{fig:stage4-network}
\end{figure}

Запустим БД и приложения в новой сети (рисунки~\ref{fig:stage4-db} и~\ref{fig:stage4-app}).

\begin{figure}[H]
    \centering
    \includegraphics[width=0.8\textwidth]{image-11.png}
    \caption{Запуск БД в новой сети}
    \label{fig:stage4-db}
\end{figure}

\begin{figure}[H]
    \centering
    \includegraphics[width=0.8\textwidth]{image-12.png}
    \caption{Запуск приложения в новой сети}
    \label{fig:stage4-app}
\end{figure}

Благодаря тому, что контейнеры запущены в одной сети, мы можем прописать адрес БД, используя имя контейнера, так как для пользовательских сетей запускается docker DNS, в отличии от default.

Проверим доступность приложения (рисунок~\ref{fig:stage4-check}).

\begin{figure}[H]
    \centering
    \includegraphics[width=0.8\textwidth]{image-13.png}
    \caption{Проверка доступности приложения}
    \label{fig:stage4-check}
\end{figure}

Докажем на практике, что контейнеры изолированы, если находятся в разных сетях, и доступны по имени, если находятся в одной сети (рисунок~\ref{fig:stage4-isolation}).

\begin{figure}[H]
    \centering
    \includegraphics[width=0.8\textwidth]{image-14.png}
    \caption{Демонстрация изоляции и доступности по имени}
    \label{fig:stage4-isolation}
\end{figure}

\section{Этап 5. Финальный артефакт}

Выведем финальный артефакт лабораторной работы (рисунок~\ref{fig:stage5-final}).

\begin{figure}[H]
    \centering
    \includegraphics[width=0.8\textwidth]{image-15.png}
    \caption{Финальный артефакт лабораторной работы}
    \label{fig:stage5-final}
\end{figure}

На рисунке~\ref{fig:stage5-final} видны:
\begin{enumerate}
    \item оба контейнера (\texttt{lyasin-postgres}, \texttt{lyasin-assistant});
    \item именованный volume (\texttt{lyasin-pgdata});
    \item пользовательская сеть (\texttt{lyasin-backend}).
\end{enumerate}

% ===== ЗАКЛЮЧЕНИЕ =====
\conclusions

В ходе практической работы была изучена экосистема Docker: освоены базовые команды, создание и использование именованных volume, bind mount для конфигурации, а также пользовательские сети для изоляции контейнеров. Продемонстрирована персистентность данных при пересоздании контейнера и работа docker DNS в пользовательской сети.

\end{document}